By making a declaration on your honour, you affirm that you have complied with
the usual rules of scientific honesty. For the formatting basically the same 
as in the previous section applies. The only addition is a field for your
signature. This template offers a separate environment for the declaration.
You can add a declaration on as shown in Listing~\ref{lst:final:statement}.

\lstinputlisting[language={[LaTeX]{TeX}},style=block,escapechar=\!,float,caption={Honourable declaration},label={lst:final:statement}]{code/final_statement.en.tex}

The example corresponds to the wording as it stood at the time of writing of
this manual footnote{April 2017} by the Central Examination Office of the
Chemnitz University of Technology. Please note that an affirmation in lieu of
an oath is is normally not necessary.

If the heading of the declaration has to be adapted, this can be done via the
optional argument of the environment. You create a signature field
by the command sequence \verb+\tucsignature+. You only have to enter your name.
For group work with several authors, you can use the command several times in
succession. If you need a wider signature line because, for example, you have a
longer name, enter a corresponding length in the optional argument, specify an
appropriate length. By default, \qty{5}{\centi\metre} is used.

The environment automatically starts a new right-hand page and also makes sure
that no page number is printed.
